\section{Related Work}
\paragraph{Coordinate-based GUI Grounding.} The core challenge for GUI agents~\citep{wang2024gui} is grounding: aligning user instructions with the correct actionable elements in the screenshot, due to the semantic inconsistency of layouts and UI elements across diverse GUI environments~\citep{liu2025infiguig1advancingguigrounding}. In early attempts, along with the screenshots, extra structured inputs, such as HTML for U-Ground~\citep{gou2024navigating}, or extractions from visual parsing modules, such as element extractions from OmniParser~\citep{Wan_2024_CVPR,yu2025omniparserv2structuredpointsofthoughtunified} and generated captions as contexts in Aria-UI~\citep{yang2024aria}, are fed into MLLMs as the supplementary inputs for the better interface understanding and more precise localization. Later works, such as AGUVIS~\citep{xu2024aguvis}, SeeClick~\citep{cheng2024seeclick} and OS-Atlas~\citep{wu2024atlas}, explore the GUI grounding leveraging MLLMs~\citep{bai2025qwen2,chen2024expanding} with screenshot-only inputs, enabling the end-to-end localization with improved scalability across diverse GUI environments. The grounding approach in these methods is to generate coordinate-based click centers or bounding boxes as text, which is indirect alignment of visual grounding requiring additional efforts, such as  scaled GUI corpora~\citep{qin2025ui,chai2024amex} and OCR pretraining~\citep{hong2024cogagent} for connecting coordinates with UI elements. Recent endeavors manage to address this gap from the data-intensive manner of the coordinate-based GUI grounding methods via incorporating GUI-specific grounding reward signals into RL-training~\citep{luo2025gui,lu2025ui,liu2025infiguig1advancingguigrounding,tang2025gui} and performing autonomous GUI exploration~\citep{fan2025gui,wu2025gui}. 

\paragraph{Coordinate-free GUI Grounding.} Recent works have begun to replace traditional coordinate‑based grounding text with patch‑level attention map predictions, where patches belonging to the correct region receive higher weight. TAG~\citep{xu2024attentiondrivenguigroundingleveraging} first leveraged the multi-head self-attention between GUI instructions and visual tokens in MLLMs for tuning‑free GUI grounding, showcasing the intrinsic potential of MLLMs' attention patterns for GUI tasks. However, the generalization of TAG for novel interfaces is restricted by the backbone and by the heuristics used for selecting text tokens and attention heads. SE-GUI~\citep{yuan2025enhancing} retains a coordinate‑based prediction manner but employs self‑attention to iteratively filter noisy training samples during training. Aside from these attention-based improvements, GUI-Actor~\citep{wu2025gui} introduces an embedding‑based grounding head that derives attention between input visual patch embeddings and the final hidden state of a special \texttt{<ACTOR>} token. Compared with \method{}, GUI-Actor adds an additional module to connect cross-layer visual and \texttt{<ACTOR>} embeddings, requiring an extra warm-up training phase and resulting in much lower training efficiency.